% chapters/intro.tex
\addcontentsline{toc}{section}{ВСТУП}
{\centering\bfseries ВСТУП\par}
\vspace{1em}

\textbf{Актуальність теми.}
Швидке перетворення Фур'є (ШПФ, англ.\ \textit{Fast Fourier Transform}) належить
до фундаментальних алгоритмів цифрового оброблення сигналів (ЦОС) і щоденно
виконується мільярди разів у мобільному зв'язку 4G/5G (OFDM-модуляція), системах
активного шумозаглушення, медичній діагностиці (МРТ, ЕЕГ-спектрометрія) та
обробленні аудіо- і відеопотоків. За оцінками IEEE Signal Processing Society
(2024), ШПФ лишається найчастіше застосовуваним числовим алгоритмом у вбудованих
системах реального часу. Алгоритм radix-4 для $N = 16$ через тристадійну
$4 \times 4$ декомпозицію потребує лише 9 нетривіальних комплексних множень
проти 256 у прямому ДПФ, що дає прискорення у 28~разів.

Попри наявність зрілих виробничих бібліотек (FFTW, ARM CMSIS-DSP, NVIDIA cuFFT,
Intel IPP), жодна з них не надає навчально-дослідницького засобу для
\emph{безпосереднього спостереження} потокової природи алгоритму: динаміки
токенів між вузлами, паралельної активації незалежних обчислювальних блоків та
наскрізної затримки у графі обчислень. Існуючі засоби відображають лише
числовий результат, приховуючи мікроархітектурний хід обчислення. Водночас
парадигма обчислень під керуванням потоку даних (dataflow computing) є природною
для ШПФ: граф метеликів безпосередньо відображається на граф потоку даних
(DFG), а чотири незалежні вузли 4-точкового ДПФ кожного ярусу активуються
паралельно без явної синхронізації. Це зумовлює актуальність розроблення
веб-орієнтованого симулятора, що поєднує достовірне виконання реального
алгоритму radix-4 з інтерактивною DFG-візуалізацією.

\textbf{Мета роботи} --- розроблення математичної моделі, графу потоку даних та
програмного симулятора 16-точкового ШПФ з основою 4 для інтерактивної
візуалізації потокового виконання алгоритму у веб-браузері.

\textbf{Об'єктом дослідження} є процеси обчислення дискретного перетворення
Фур'є під керуванням потоку даних у програмних симуляторах обчислювальних систем.

\textbf{Предметом дослідження} є методи та підходи до математичного моделювання
алгоритму radix-4, побудови графу потоку даних (DFG / SDF) обчислювача та
програмної реалізації інтерактивного симулятора на платформі веб-браузера.

\textbf{Задачі дослідження:}
\begin{enumerate}
  \item проаналізувати сімейство алгоритмів ШПФ, парадигму dataflow computing та
    методи візуального моделювання обчислювальних систем; виконати системний
    аналіз предметної області та обґрунтувати вибір архітектурного підходу;
  \item розробити математичну модель алгоритму radix-4 для $N = 16$ з
    відображенням індексів, трьома стадіями обчислення, матрицею поворотних
    множників та оцінкою обчислювальної складності;
  \item розробити граф потоку даних обчислювача зі специфікацією вузлів, дуг,
    затримок і керуючих слів; визначити критичний шлях та коефіцієнт
    паралелізму DFG;
  \item реалізувати ядро симуляції \texttt{DataflowRuntime} з правилом активації
    токенів та веб-інтерфейс візуалізації на React Flow з анімацією потоку
    токенів і тепловою картою активності;
  \item виконати верифікацію результатів симулятора порівнянням із еталонним ДПФ
    бібліотеки \texttt{fft.js} із точністю $\varepsilon = 10^{-9}$ на шести
    тестових сценаріях;
  \item провести експериментальний аналіз обчислювальної складності та
    паралелізму dataflow-моделі; підтвердити теоретичні оцінки рантаймовими
    метриками симулятора.
\end{enumerate}

\textbf{Методологія та інструментарій.} Математична база дослідження --- теорія
дискретних перетворень Фур'є та алгоритмів Cooley--Tukey, формалізм синхронних
dataflow-графів (SDF) Е.~Лі і Е.~Мессершмітта. Інженерна основа --- архітектурне
моделювання за нотацією C4 та UML-діаграмами класів і послідовностей. Програмна
реалізація використовує стек Next.js~14 / React~18 / TypeScript~5 із бібліотеками
React Flow (\texttt{@xyflow/react}) 12 для рендерингу графу, Zustand~4 для
реактивного управління станом, \texttt{fft.js} як еталон верифікації та Jest~29
для модульного й інтеграційного тестування.

\textbf{Практична цінність роботи.} Розроблений симулятор є
навчально-дослідницьким інструментом для вивчення потокових алгоритмів ЦОС у
курсі «Цифрове оброблення сигналів» та суміжних дисциплінах бакалаврської
програми 122 «Комп'ютерні науки». Застосунок виконує реальний алгоритм radix-4
(не його імітацію), доступний без встановлення додаткового програмного
забезпечення через веб-браузер і підтримує регульовану швидкість симуляції та
покроковий режим для прецизійного дослідження активацій вузлів. Верифікована
точність $\varepsilon < 10^{-9}$ підтверджує придатність для прикладних числових
досліджень. Архітектура DFG (56 вузлів, 64 дуги, параметри \texttt{latency} і
\texttt{delay}) може слугувати вихідною специфікацією для апаратної реалізації
обчислювача на FPGA, де граф потоку даних безпосередньо відображається на
маршрутизацію сигналів між DSP-блоками.

\textbf{Апробація результатів.} Окремих публікацій за темою роботи не подавалось;
результати роботи планується подати на студентську науково-технічну конференцію
НУ «Львівська політехніка» у 2026~році.